\begin{tomtat}
\begin{itemize}
  \item Định nghĩa~\ref{def:ch08-chi-so-trung-thuc} gọi
  \tn{chỉ số trung thực}{faithfulness metric} là một thủ tục trả về một số kèm
  quy ước đọc số ấy là mức độ trung thực. Định nghĩa
  không đòi bằng chứng nào cho quy ước, và chương này đi tìm xem từng họ chỉ số
  lấp chỗ hở đó bằng gì. Mọi chỉ số ở đây thuộc ô đánh giá theo chức năng của
  Doshi-Velez và Kim, mà hai tác giả nói là thích hợp nhất khi đối tượng đem đo
  đã được kiểm chứng từ trước bằng thí nghiệm có người tham gia.
  \item Hình~\ref{fig:ch08-vong-lap-do} rút mọi chỉ số về một vòng lặp với ba
  lựa chọn tự do: chọn tập con nào theo thứ tự điểm, thay các đặc trưng bị bỏ
  bằng giá trị gì, và so đại lượng nào. Các họ chỉ số khác nhau đúng ở ba chỗ
  ấy.
  \item Họ xóa dần gồm deletion và insertion của RISE, comprehensiveness và
  sufficiency của ERASER. Bốn chỉ số ấy rơi vào hai đại lượng mà
  chương~\ref{ch:khung-hop-nhat} đã nhận: insertion và sufficiency giữ lại phần
  điểm cao nên đo fidelity, còn deletion và comprehensiveness bỏ phần điểm cao
  nên đo inverse fidelity. RISE chọn ảnh làm mờ thay cho ảnh
  trống, và ERASER tự ghi rằng đây là những chỉ số đơn giản cho một câu hỏi còn
  mở.
  \item ROAR chỉ ra rằng đầu vào đã bỏ bớt đặc trưng nằm ngoài phân phối huấn
  luyện, nên không phân biệt được hiệu năng giảm vì mất thông tin hay vì dấu vết
  lạ, và nó huấn luyện lại mô hình để gỡ chỗ lẫn ấy. Cái giá là đại lượng được
  đo đổi sang một mô hình khác. Hase và cộng sự siết thêm: tiên nghiệm của mô
  hình và cách khởi tạo trọng số lọt được vào chính các chỉ số.
  \item Sensitivity-n và infidelity ở phương trình~\ref{eq:ch08-infidelity} đòi
  một đẳng thức trên nhiều tập con thay vì một thứ tự, nhưng vẫn kiểm nó qua
  phản ứng của mô hình trước phép nhiễu, và phân phối nhiễu vẫn do người dùng
  chọn. Hai \tn{phép thử ngẫu nhiên hóa}{randomization test} của Adebayo và cộng
  sự là điều kiện cần chứ không phải thang đo: trượt thì bị loại, qua thì chưa kết luận gì.
  \item Bài báo neo chia các chỉ số thành bốn nhóm theo chỗ chúng can thiệp, và
  chẩn đoán rằng nhóm dựa trên mức quan trọng đồng nhất mức quan trọng với độ
  trung thực, còn nhóm dựa trên \tn{tính hữu dụng ngữ nghĩa}{semantic utility}
  bị một lời biện minh bịa đặt qua mặt. Đối tượng của bài báo ấy là chuỗi suy luận chứ không phải
  attribution đặc trưng, nên cái đi qua được ranh giới là cấu trúc của lập luận,
  không phải kết quả đo.
\end{itemize}
\end{tomtat}

\cauhoi
\begin{cauhoids}
  \item Với mỗi họ chỉ số trong chương, hãy điền ba ô viền đứt của
  hình~\ref{fig:ch08-vong-lap-do}: tập con được chọn thế nào, giá trị thay thế
  là gì, và đại lượng nào được so. Hai họ nào trùng nhau ở hai ô và chỉ khác một
  ô?
  \item Deletion và comprehensiveness được đề xuất ở hai lĩnh vực khác nhau, cho
  hai loại dữ liệu khác nhau, dưới hai cái tên không gợi gì đến nhau. Hãy lập
  luận vì sao chúng vẫn đo cùng một thứ, và nêu một trường hợp mà hai chỉ số ấy
  cho kết luận ngược nhau.
  \item ROAR huấn luyện lại mô hình trước khi đo. Hãy chỉ ra chỗ mà việc huấn
  luyện lại làm đại lượng được đo lệch khỏi
  định nghĩa~\ref{def:ch01-do-trung-thuc}, rồi cho biết vì sao lệch như vậy vẫn
  có thể là đánh đổi đáng giá.
  \item Một phương pháp attribution đạt infidelity rất thấp theo
  phương trình~\ref{eq:ch08-infidelity} với một phân phối $\mu_I$ do chính người
  đề xuất phương pháp ấy chọn. Hãy nêu điều kết quả đó chứng minh và điều nó
  không chứng minh.
  \item Hai phép thử ngẫu nhiên hóa là điều kiện cần. Hãy dựng một ví dụ về một
  phương pháp qua cả hai phép thử mà vẫn sai hoàn toàn về cơ chế, và cho biết
  loại bằng chứng nào mới loại được nó.
  \item Đọc \enquote{A Benchmark for Interpretability Methods in Deep Neural
  Networks} của Hooker và cộng sự~\autocite{mroar}. Các tác giả nêu lý do nào
  cho việc phải huấn luyện lại, và lý do ấy có áp dụng được cho một mô hình mà
  ta không có quyền huấn luyện lại hay không?
  \item Đọc \enquote{Faithfulness Metrics Don't Measure Faithfulness: A
  Meta-Evaluation with Ground Truth} của Gur-Arieh và cộng
  sự~\autocite{p21metrics}, phần mô tả bốn nhóm chỉ số. Hãy xếp bốn họ trong
  chương này vào bốn nhóm ấy, chỉ ra nhóm nào của bài báo không có đại diện
  trong chương, và giải thích vì sao đối tượng khác nhau lại dẫn tới chỗ trống
  đó.
\end{cauhoids}
